\uuid{nAee}
\titre{Dénombrement}
\auteur{bodin ; hanani ; mzari}
\organisation{exo7}
\licence{CC BY-NC-SA 4.0}
\datecreate{2020-01-20}
\niveau{L1}
\module{Algèbre}
\chapitre{Logique, ensemble, raisonnement}
\sousChapitre{Ensemble}
\difficulte{1}
\theme{Ensembles, applications, dénombrement}
\exo7qcmid{lille-1.ensembles.008}

\contenu{
\begin{question}
\qclassification{100.02,101.01,102.01,102.02}
\qtags{motcle=dénombrement}
On désigne par $\mathrm{C}^k_n$ le nombre de choix de $k$ éléments parmi $n$. Combien fait $\displaystyle \sum _{k=0}^{10}\mathrm{C}^k_{10}$ ?
\begin{answers}
    \bad{$10$}
    \bad{$100$}
    \good{$1024$}
    \bad{$50$}
\end{answers}
\begin{explanations}
Le binôme de Newton donne $\displaystyle \sum _{k=0}^{10}\mathrm{C}^k_{10}=(1+1)^{10}=2^{10}=1024$.
\end{explanations}
\end{question}
}
